\documentclass{article}
\usepackage[a4paper]{geometry}
\usepackage{fancyhdr}
\pagestyle{fancy}
\lhead{Die Abiturprüfung}
\rhead{April 2026}

\usepackage{hyperref}
\begin{document}
\section{Die Abiturprüfung}
\subsection{Auswahl}
Es wird eine Klausur gewählt, dessen Themen einfach sind und dessen dessen Themenübergriff bekannt ist.
 
\subsection{Aufgabe 3}
Die Qualität der Argumentation ist hier erstmal wichtiger als die Quantität.

Es werden Kritierien festgelegt, je nach Operator, und anhand für alle Optionen Pro- und Contra-Argumente gefunden.

Als Fazit wird ein Urteil gebildet.

\subsection{Themen}
\begin{description}
 \item[Politische Partizipation]~
 \begin{itemize}
   \item \hyperref[Partizipationsformen Analysieren]{Funktionen}
   \item Formen
 \end{itemize}
 \item[\mbox{\hyperref[Soziale Marktwirtschaft]{Soziale Marktwirtschaft}}]~
 \begin{itemize}
   \item Prinzipien
   \item Aufgaben des Staates
   \item Gerechtigkeit
   \item Umwelt
 \end{itemize}
 \item[Friedenssicherung]~
 \begin{itemize}
   \item \hyperref[Die Bundeswehr]{Auslandseinsätze}
   \item Verfassungsrecht
   \item \hyperref[Alte und Neue Kriege]{Alte und Neue Kriege}
   \item \hyperref[Das Zivilisatorische Hexagon]{Zivilisatorisches Hexagon}
   \item \hyperref[Fragile Staatlichkeit]{Fragile Staatlichkeit}
   \item Bündnisverpflichtungen 
 \end{itemize}
 \item[Wirtschaftliche Globalisierung]~
 \begin{itemize}
   \item \hyperref[Freihandel und Protektionismus]{Freihandel und Protektionismus}
   \item \hyperref[Handelsregime]{Handelsregime}
 \end{itemize}
\end{description}

\end{document}